% Member 3 - final report, IV.E.1 Failure Cases
% The numbers are read from figures/failure_cases/index.json, which is
% written by member4_training_eval/make_failure_cases.py.
\subsection{Failure Cases}
\label{sec:failure_cases}

The mechanism analysis of Section~\ref{sec:mechanism} established that the
spatial modulation does not preserve structure selectively. This section
shows what that failure looks like on individual images, because the
aggregate statistics hide the shape of it.

\paragraph{How the cases are chosen.}
The twelve images of Fig.~\ref{fig:failure_cases} are not a hand-picked
selection. They are the twelve images, out of all 1{,}349, with the
smallest paired EdgeScore difference between the spatial variant and the
strength-matched uniform control: the cases on which the modulation does
the most damage. The selection uses only image-level measurements, never a
target-domain accuracy, so no evaluation image influences it.

\paragraph{What the panels show.}
Each row of Fig.~\ref{fig:failure_cases} gives, from left to right, the
content image, the style reference, the structure map $M$, the strength
field $\alpha$, the uniform AdaIN result and the spatial result. The
pattern is consistent across the rows, and it is not the pattern the method
was designed to produce.

The structure map responds strongly to the \emph{background} of these
images: the foliage, the brickwork and the canvas texture of the paintings
all produce large responses, while the object's own contour often does not
stand out above them. Because $\alpha = \alpha_0(1-\beta M)$ treats every
response alike, the modulation spends its protection budget on the
background. The alpha map makes this visible directly: its dark bands
follow the background texture rather than the outline of the object.

The result is that the object's shape receives the full strength
$\alpha_0 = 0.5$ in the places where it mattered most, while the
background, which needed no protection, receives less. The matched
uniform control, applying the same strength everywhere at the same mean,
keeps more of the object's contour simply because it applies it
uniformly. That is why the spatial variant's EdgeScore is \emph{lower}
than the matched control's on these images, and it is why the aggregate
difference of Section~\ref{sec:mechanism} is negative rather than merely
insignificant.

\paragraph{Where the failures concentrate.}
The class distribution of the twelve worst cases is itself informative:
five are \texttt{person} and four are \texttt{guitar}, the two classes with
the lowest edge density in PACS. A \texttt{person} photograph is dominated
by a textured background and a loosely bounded silhouette; a \texttt{guitar}
photograph is dominated by a long, thin neck against a similarly textured
background. Both are exactly the situations in which an edge response
fails to separate the class-relevant contour from the rest of the image.
The classes that the modulation handles best are the ones with a compact,
high-contrast object and a plain background, which is the easy case for any
edge-based prior.

\paragraph{What the cases imply.}
The failure is not a bug in the implementation: the $\beta = 0$ cell
reproduces plain AdaIN byte for byte, so the generator does what it is
told. It is a limitation of the \emph{prior}. A soft Sobel response is
not a foreground mask, and the method's premise --- that a response
magnitude identifies the structure worth protecting --- does not hold for
the images where protection would matter most. Section~\ref{sec:negative_results}
draws the consequence for the classification results; this section shows
the mechanism behind it.
